\begin{resumo}
A terceirização do desenvolvimento de \textit{software} na Administração Pública brasileira é amplamente adotada como estratégia de otimização de custos e transferência da carga operacional para a iniciativa privada. Contudo, a delegação do desenvolvimento não isenta o poder público da responsabilidade sobre a qualidade do produto entregue, cujas deficiências impactam diretamente o cidadão. O normativo vigente, composto pela Lei nº 14.133/2021 e pela Portaria SGD/MGI nº 750/2023, estabelece diretrizes para essas contratações, mas delega ao gestor do contrato a definição dos modelos, métricas e critérios de aferição da qualidade — configurando uma lacuna operacional crítica para a fiscalização técnica de entregas. Nesse contexto, este trabalho investiga de forma empírica como operacionalizar os atributos de qualidade de Manutenibilidade, em especial Modularidade e Testabilidade, por meio de indicadores objetivos capazes de apoiar a aceitação de versões de produtos de \textit{software} desenvolvidos por terceiros. Para tanto, foi conduzida uma Revisão Estruturada da Literatura para insumos bibliográficos. E a definição do planejemanto do Estudo de caso para a avaliação quantitativa  por meio da extração automatizada de métricas do histórico de \textit{releases} integradas ao repositório, agregadas segundo os preceitos do meta-modelo MeasureSoftGram.

 \vspace{\onelineskip}

 \noindent
 \textbf{Palavras-chave}: qualidade de software. manutenibilidade. terceirização. contratos públicos. métricas de código. análise estática. estudo de caso.
\end{resumo}
